\section{Logic}

\subsection{What is an argument?}

\begin{readingnotes}
	"Having an argument" is different from a logical argument.
\end{readingnotes}

\begin{keyquote}
	Sally: Abortion is morally wrong because it is wrong to take the life of an
	innocent human being, and a fetus is an innocent human being.
\end{keyquote}

\begin{readingnotes}
	This shows an actual logical argument. The \textbf{conclusion} of the argument
	is the first four words "abortion is morally wrong".

	The reason for thinking the conclusion is true is called the \textbf{premise}.

	Both premises and conclusions are statements. A \textbf{statement} is a type of
	sentence that can be true or false, and is a declarative sentence.
	ex: The Nile is a river in northeast Africa.\\

	Sally's argument can be captured as follows:
	\begin{enumerate}
		\item It is morally wrong to take the life of an innocent human being
		\item A fetus is an innocent human being
		\item Therefore, abortion is morally wrong
	\end{enumerate}

	By convention, the last numbered statement should be the conclusion. This form
	is called \textbf{standard argument form}.

	An \textbf{argument} is a set of statements where the premises attempt to
	provide a reason for thinking that the conclusion is true.
\end{readingnotes}

\subsection{Identifying Arguments}

\begin{readingnotes}
	Best way to identify an argument is to ask if there is a statement established
	as true by basing it on some other statement. This means, generally, that
	arguments have to be 2+ statements.

	Look for premise indicators like "because" or "since", and conclusion indicators
	like "so" or "therefore". Indicators are not failsafe though.
\end{readingnotes}

\begin{table}[htb]
	\centering
	\caption{List of common premise and conclusion indicators}
	\begin{tabular}{|l|l|}
		\hline
		\textbf{Premise indicators} & \textbf{Conclusion indicators} \\ \hline
		since                       & therefore                      \\ \hline
		because                     & so                             \\ \hline
		for                         & hence                          \\ \hline
		as                          & thus                           \\ \hline
		given that                  & implies that                   \\ \hline
		seeing that                 & consequently                   \\ \hline
		for the reason that         & it follows that                \\ \hline
		is shown by the fact that   & we may conclude that           \\ \hline
	\end{tabular}
	\label{tab:indicators}
\end{table}

\begin{readingnotes}
	The \textbf{substitution test} is a way to figure out if a word is a premise/conclusion
	indicator. Just substitute one of the indicators from table \ref{tab:indicators},
	and see if the resulting sentence still makes sense.
\end{readingnotes}

\subsection{Validity}

\begin{readingnotes}
	A \textbf{valid argument} is an argument whose conclusion cannot be false,
	assuming all of the premises are true. To test validity, take all premises as
	true, then see if the conclusion is true. If it is, then the argument is valid,
	if not, then it is not valid. This has nothing to do with the actual truth
	of the premises or argument in general (that comes up with soundness).

	A \textbf{counterexample} is a description of a scenario where the premises of
	the argument are all true, while the conclusion of the argument is false.
	Counterexamples automatically invalidate arguments.

	The \textbf{informal test of validity} is a test where you try to imagine
	a world where all of the premises are true and yet the conclusion is false.
	If you can imagine this world, then the argument is invalid.
\end{readingnotes}
